\documentclass[12pt]{article}
\usepackage[margin=1in]{geometry}
\usepackage{amsmath, amssymb, amsthm}
\usepackage{booktabs}
\usepackage{hyperref}
\usepackage{graphicx}
\usepackage{array}
\usepackage{xcolor}
\usepackage{titlesec}
\usepackage{abstract}
\usepackage{multirow}
\usepackage{float}
\usepackage{tikz}
\usetikzlibrary{arrows.meta, positioning, shapes}
\usepackage{listings}
\usepackage{enumerate}

\hypersetup{
    colorlinks=true,
    linkcolor=black,
    urlcolor=blue,
    citecolor=black
}

% ── Code listing style ──────────────────────────────────────────────────────
\definecolor{codebg}{RGB}{15,15,15}
\definecolor{codecomment}{RGB}{100,160,100}
\definecolor{codestring}{RGB}{200,150,80}
\definecolor{codekw}{RGB}{100,180,255}
\definecolor{codenumber}{RGB}{160,160,160}

\lstset{
  language=Python,
  backgroundcolor=\color{codebg},
  basicstyle=\ttfamily\footnotesize\color{white},
  commentstyle=\color{codecomment},
  stringstyle=\color{codestring},
  keywordstyle=\color{codekw},
  numberstyle=\tiny\color{codenumber},
  numbers=left,
  numbersep=8pt,
  breaklines=true,
  frame=single,
  framerule=0.4pt,
  rulecolor=\color{gray},
  xleftmargin=1.5em,
  framexleftmargin=1.5em,
  showstringspaces=false,
  tabsize=4,
  captionpos=b,
}

\newtheorem{theorem}{Theorem}
\newtheorem{definition}{Definition}
\newtheorem{proposition}{Proposition}

% ── Title ───────────────────────────────────────────────────────────────────
\title{\textbf{Building SPECTRUM}\\
\large Design Rationale and Implementation of a\\
Multi-Substrate AI Architecture with Neuromorphic Components}
\author{
    Winston Zai Lin \\
    \small Blair Academy \\
    \small \texttt{linzaiwinston413@gmail.com}
}
\date{April 2026}

\begin{document}
\maketitle

% ── Abstract ────────────────────────────────────────────────────────────────
\begin{abstract}
This paper has two goals: to explain \textit{why} the SPECTRUM multi-substrate AI architecture is designed the way it is, and to show \textit{how to build} its core neuromorphic component. The first half derives the SPECTRUM architecture from first principles---the substrate problem, why transformers are the wrong tool for reasoning, and what a physically heterogeneous system looks like. The second half is an implementation guide. We present the DST (Dendritic/Somatic/Threshold) neuron, a three-variable model designed from architectural principles rather than biological replication, and its extension to DSTC (adding a two-compartment dendrite, voltage-gated NMDA coincidence detection, inhibitory synapses, and a calcium plasticity signal). We provide complete Python implementations with annotated equations, working simulations, and the full circuit specification for physical fabrication. The D compartment---an electrolytic capacitor and bleeder resistor---is the first buildable piece; its schematic, component values, and test protocol are given explicitly. SPECTRUM's position is that the ceiling of AI capability is at least partly a substrate problem rather than a purely algorithmic one; this paper does not attempt to establish that, and states it as the motivation for the design rather than as a result. What the paper does establish is narrower and measured: against LIF, Izhikevich and AdEx under matched conditions, DST is the only model tested that holds an absolute refractory period, and its linear dynamics tolerate a coarser integration timestep and lower cost than either nonlinear competitor---which is also what makes the D compartment realisable as a capacitor and a resistor.
\end{abstract}

% ── 1. Why SPECTRUM Exists ──────────────────────────────────────────────────
\section{Why SPECTRUM Exists: The Substrate Problem}
\label{sec:why}

Every major AI system today runs on the same substrate: digital silicon executing floating-point matrix multiplication. This is not a principled design choice. It is an accident of history. The transformer architecture \cite{vaswani2017} achieved dominance because it maps cleanly onto GPU hardware---attention is a large matrix multiply, and GPUs are optimized for exactly that operation.

The limitation is invisible until you look for it. Consider what intelligence actually needs to do:

\begin{itemize}
    \item \textbf{Detect events} in a stream of sensory input---spikes arriving milliseconds apart.
    \item \textbf{Retrieve memories} from partial or corrupted cues---pattern completion under noise.
    \item \textbf{Reason continuously}---run differential equations, simulate physical dynamics, optimize over smooth spaces.
    \item \textbf{Search} large hypothesis spaces---combinatorial optimization.
    \item \textbf{Store} long-term knowledge---persist information over days, months, years.
    \item \textbf{Communicate}---translate internal reasoning into human language.
\end{itemize}

A digital transformer does all six with the same operation: matrix multiplication over 32-bit floats. It does some of them tolerably. It does none of them optimally. Each of those six functions has a physical substrate for which it is the natural computation:

\begin{center}
\small
\begin{tabular}{ll}
\toprule
\textbf{Function} & \textbf{Optimal substrate} \\
\midrule
Event detection & Neuromorphic (spike-triggered, not clock-triggered) \\
Pattern completion & Energy-based / Hopfield networks \\
Continuous reasoning & Analog circuits (differential equations in physics) \\
Search & Quantum (Grover's $\sqrt{N}$ speedup) \\
Long-term storage & Molecular (DNA, near-zero passive energy) \\
Language output & Digital transformer \\
\bottomrule
\end{tabular}
\end{center}

\noindent\textbf{SPECTRUM's central claim}: build a system that routes each function to its physically appropriate substrate. The transformer becomes the output interface---the mouth---not the brain.

\subsection{The Energy Argument}

In deep-reasoning mode (an autonomous agent running $K$ internal cycles before producing one language output), the energy comparison is stark. A transformer-only system runs $K$ full forward passes at $\sim$10W each. SPECTRUM runs $K$ cycles on analog ($\sim$5mW) and Hopfield ($\sim$50mW) substrates, invoking the transformer only once at the end:

\[
E_{\text{transformer}} = 10\,\text{W} \cdot K, \qquad
E_{\text{SPECTRUM}} \approx 55\,\text{mW} \cdot K + 10\,\text{W}
\]

\noindent The asymptotic ratio is $\approx 180\times$ at current analog hardware, scaling to $10^3$--$10^5\times$ for mature analog substrates. The advantage does not exist for single-query chat; it is decisive for long-horizon autonomous reasoning.

% ── 2. The SPECTRUM Architecture ────────────────────────────────────────────
\section{The SPECTRUM Architecture}
\label{sec:spectrum}

SPECTRUM (Synaptic Parallel Ecosystem with Causal Temporal Reasoning Unit Mesh) defines six layers, each on a physically distinct substrate:

\begin{table}[H]
\centering
\small
\begin{tabular}{lllll}
\toprule
\textbf{Layer} & \textbf{Substrate} & \textbf{Time scale} & \textbf{Power} & \textbf{Function} \\
\midrule
L1: Perception     & Neuromorphic (DSTC) & $\mu$s--ms  & $<$10mW  & Event detection, coincidence \\
L2: Working Memory & Hopfield networks   & ms--s       & $\sim$50mW & Pattern completion, recall \\
L3: Long Memory    & Molecular (DNA)     & s--years    & $\sim$0   & Persistent knowledge \\
L4: Reasoning      & Analog              & s--min      & $\sim$5mW  & Continuous optimization \\
L5: Search         & Quantum             & ms--s       & variable  & Hypothesis space search \\
L6: Language       & Digital (ARBOR)     & s           & $\sim$10W  & Output, routing, synthesis \\
\bottomrule
\end{tabular}
\caption{SPECTRUM layer specification. All layers communicate through a shared 1024-dimensional continuous embedding space. The only cross-layer coupling is this embedding; each substrate is otherwise physically independent.}
\label{tab:layers}
\end{table}

\noindent Three design principles govern every choice:

\begin{enumerate}
    \item \textbf{Language is the output, not the reasoning.} L6 is the interface. The actual computation happens in L1--L5.
    \item \textbf{Time scales are explicit.} Each layer operates at its natural timescale. Nothing is forced to a single clock.
    \item \textbf{Memory and compute are substrate-separated.} L3 stores; L2 retrieves; L4 reasons. No layer does all three.
\end{enumerate}

\noindent\textbf{Information flow.} Input enters at L1 (fast perceptual processing). Salient patterns pass to L2 (working memory). L2 retrieves relevant context from L3. The assembled context enters L4 (reasoning), which may invoke L5 (search) for optimization subproblems. L6 translates the result into language. Critically, L4 can loop back to L1, L2, and L3 internally without producing a language output---internal reasoning is substrate-efficient.

% ── 3. The DST Neuron: Design from Scratch ──────────────────────────────────
\section{The DST Neuron: Design from Scratch}
\label{sec:dst}

L1 requires a neuromorphic unit: a circuit that fires spikes in response to input, integrates inputs over time, and resets after firing. The standard approach borrows from biology---the Hodgkin-Huxley (HH) model \cite{hodgkin1952}, with four coupled variables and six ionic conductance functions calibrated to squid axon data. HH is accurate and biologically faithful. It is also the wrong starting point for a chip.

We derive the DST neuron differently: from the three computational jobs a neuron has to do.

\subsection{Three Jobs, Three Variables}

\begin{description}
    \item[$D(t)$ --- Dendritic potential.] The dendrite integrates incoming synaptic currents over a time window. Inputs that arrive close together should add up; inputs spread far apart should not. This is leaky integration:
    \[
    \tau_D \frac{dD}{dt} = -D + I_{\text{syn}}(t)
    \]
    $D$ charges toward $I_{\text{syn}}$ and leaks back to zero with time constant $\tau_D = 10$\,ms.

    \item[$S(t)$ --- Somatic potential.] The soma decides whether accumulated dendritic drive is sufficient to fire. It sits at rest ($S_{\text{rest}} = -1$) and is pushed toward threshold by $D$, resisted by the adaptive threshold $\theta$:
    \[
    \tau_S \frac{dS}{dt} = -(S - S_{\text{rest}}) + \kappa D - (\theta - \theta_0)
    \]
    Fire condition: $S(t) \geq 0$. After firing: $S \leftarrow S_{\text{reset}} = -2$.

    \item[$\theta(t)$ --- Adaptive threshold.] After each spike, $\theta$ jumps by $\Delta\theta = 6$, making it temporarily impossible to fire again (absolute refractory). $\theta$ then decays back to baseline $\theta_0 = 1$ with time constant $\tau_\theta = 80$\,ms (relative refractory):
    \[
    \frac{d\theta}{dt} = -\frac{\theta - \theta_0}{\tau_\theta}, \quad \theta(t_k^+) = \theta(t_k^-) + \Delta\theta
    \]
\end{description}

\noindent Together these three equations produce: single action potentials, temporal summation, spatial summation, absolute and relative refractory periods, and spike-rate adaptation---the full behavioral repertoire of a spiking neuron, with three variables instead of HH's four plus six conductance functions.

\subsection{Analytical Spike Waveform}

HH simulates the spike shape by integrating ion channel dynamics at every timestep. DST takes a different approach: the spike waveform is \textit{analytical}. At each firing time $t_k$, the displayed voltage is:

\[
V(t) = V_{\text{rest}} + A_{\text{peak}} \cdot \mathcal{G}(t - t_k,\, \sigma_{\text{rise}}) - A_{\text{AHP}} \cdot \mathcal{G}(t - t_k - \tau_{\text{AHP}},\, \sigma_{\text{AHP}})
\]

\noindent where $\mathcal{G}(x, \sigma) = \exp(-x^2 / 2\sigma^2)$. Two Gaussians: one for the depolarization peak ($A_{\text{peak}} = 100$\,mV, $\sigma_{\text{rise}} = 0.25$\,ms), one for the afterhyperpolarization ($A_{\text{AHP}} = 12$\,mV, $\tau_{\text{AHP}} = 1.2$\,ms offset, $\sigma_{\text{AHP}} = 2$\,ms). This waveform is computed in $O(N_{\text{fires}})$ rather than being integrated at every timestep. It is $\sim$10$\times$ faster than HH and produces a cleaner, tunable spike shape.

\subsection{Parameters}

\begin{center}
\small
\begin{tabular}{lll}
\toprule
\textbf{Parameter} & \textbf{Value} & \textbf{Meaning} \\
\midrule
$\tau_D$    & 10 ms  & Dendritic integration window (EPSP decay) \\
$\tau_S$    &  2 ms  & Somatic response speed \\
$\tau_\theta$ & 80 ms & Threshold recovery (refractory timescale) \\
$\kappa$    &  1.2   & Dendritic-to-somatic coupling gain \\
$S_{\text{rest}}$ & $-1.0$ & Somatic resting level \\
$\theta_0$  &  1.0   & Baseline threshold \\
$\Delta\theta$ & 6.0 & Threshold jump per spike \\
$S_{\text{reset}}$ & $-2.0$ & Somatic reset after fire \\
\bottomrule
\end{tabular}
\end{center}

% ── 4. The DSTC Extension ────────────────────────────────────────────────────
\section{The DSTC Extension: Four Additions}
\label{sec:dstc}

DST is minimal and fast. DSTC adds four mechanisms that make the chip simulate real neurons more faithfully---and, more importantly, give L1 computational capabilities beyond simple threshold integration.

\subsection{Two-Compartment Dendrite}

Real dendrites are not a single point. Distal inputs (far from the soma) integrate more slowly than proximal inputs (near the soma). We split $D$ into two compartments in cascade:

\begin{align}
\tau_{D_d} \frac{dD_d}{dt} &= -D_d + I_{\text{AMPA}}(t) + g_{\text{NMDA}}(D_d)\,I_{\text{NMDA}}(t) - I_{\text{inh},d}(t) \label{eq:Dd} \\
\tau_{D_p} \frac{dD_p}{dt} &= -D_p + I_{\text{prox}}(t) + \kappa_{dd}\,D_d - I_{\text{inh},p}(t) \label{eq:Dp}
\end{align}

\noindent with $\tau_{D_d} = 15$\,ms (distal, slower) and $\tau_{D_p} = 5$\,ms (proximal, faster). Distal signals travel through $D_d \to D_p$ before reaching $S$, introducing a natural propagation delay. Proximal inhibition ($I_{\text{inh},p}$) is more effective than distal inhibition ($I_{\text{inh},d}$) because it intercepts the signal closer to the soma---consistent with known shunting inhibition biology.

\subsection{NMDA-Like Voltage Gate}

The NMDA receptor is biologically unique: it requires both a presynaptic spike \textit{and} postsynaptic depolarization to open. It is a coincidence detector. We implement the Mg$^{2+}$ block relief as:

\[
g_{\text{NMDA}}(V) = \frac{1}{1 + \gamma \cdot e^{-\beta V}}, \quad \gamma = 3.0,\; \beta = 2.5
\]

At rest ($V_{D_d} = -1$): $g_{\text{NMDA}} \approx 0.027$---nearly blocked. At depolarization ($V_{D_d} = 2$): $g_{\text{NMDA}} \approx 0.98$---fully open. An NMDA-only input arriving at a resting dendrite produces almost no effect. The same input arriving when the dendrite is already depolarized (e.g., by concurrent AMPA input) fires the neuron. This is hardware-level coincidence detection: L1 detects correlated inputs without digital computation.

\textbf{Hardware implementation}: a MOSFET whose gate voltage is $V_{D_d}$. The source-drain conductance varies with gate voltage exactly as $g_{\text{NMDA}}(V)$ varies with $V_{D_d}$.

\subsection{Inhibitory Synapses}

Inhibitory inputs pull the dendritic potential toward a reversal potential $E_{\text{inh}} = -2.5$ (below $S_{\text{rest}}$):

\[
I_{\text{inh}}(t) = g_{\text{inh}}(t) \cdot (V - E_{\text{inh}})
\]

This is conductance-based (shunting) inhibition: the inhibitory effect is strongest when the membrane is depolarized and weakest at rest---mirroring biological GABA$_\text{A}$ receptor dynamics. Hardware: NPN transistor current sink.

\subsection{Calcium Compartment and STDP}

A calcium compartment $Ca(t)$ rises by $\delta_{\text{Ca}} = 1$ at each spike and decays with $\tau_{\text{Ca}} = 200$\,ms:

\[
\frac{dCa}{dt} = -\frac{Ca}{\tau_{\text{Ca}}}, \quad Ca(t_k^+) = Ca(t_k^-) + \delta_{\text{Ca}}
\]

$Ca$ tracks recent spike history. It is the plasticity signal for spike-timing-dependent plasticity (STDP): the weight of a synapse increases if the presynaptic spike arrives before the postsynaptic spike (LTP) and decreases otherwise (LTD), with an exponential time window:

\[
\Delta w = \begin{cases}
A_+ \cdot e^{-\Delta t / \tau_+} & \Delta t > 0 \;\text{(LTP)} \\
-A_- \cdot e^{\Delta t / \tau_-} & \Delta t < 0 \;\text{(LTD)}
\end{cases}
\]

\noindent This enables L1 to \textit{learn online}---adapting synaptic weights during inference without invoking the digital transformer at L6. Hardware: a slow RC circuit (200\,$\mu$F $\times$ 1\,k$\Omega$, $\tau = 200$\,ms) charged by the spike output of the $S$ compartment.

\subsection{Full DSTC System}

The complete five-variable system is:

\begin{align}
\tau_{D_d} \dot{D}_d &= -D_d + I_{\text{AMPA}} + g_{\text{NMDA}}(D_d) I_{\text{NMDA}} - g_{\text{inh},d}(D_d - E_{\text{inh}}) \nonumber \\
\tau_{D_p} \dot{D}_p &= -D_p + I_{\text{prox}} + \kappa_{dd} D_d - g_{\text{inh},p}(D_p - E_{\text{inh}}) \nonumber \\
\tau_S \dot{S} &= -(S - S_{\text{rest}}) + \kappa D_p - (\theta - \theta_0) \nonumber \\
\dot{\theta} &= -(\theta - \theta_0)/\tau_\theta \nonumber \\
\dot{Ca} &= -Ca/\tau_{Ca} \nonumber
\end{align}

\noindent with spike events at $S \geq 0$: $S \leftarrow S_{\text{reset}}$, $\theta \leftarrow \theta + \Delta\theta$, $Ca \leftarrow Ca + \delta_{Ca}$.

% ── 5. Python Implementation ─────────────────────────────────────────────────
\section{Python Implementation}
\label{sec:code}

The full implementation lives in three files. We present the key functions here.

\subsection{\texttt{dst\_model.py}: The Base Simulator}

The core simulator integrates the three DST equations with Euler's method. Synaptic inputs are precomputed as exponentially-decaying currents from a list of spike times:

\begin{lstlisting}[caption={DST simulator core (\texttt{dst\_model.py})}, label={lst:dst}]
def simulate(I_func, synapses=None, t_end=100, dt=0.01):
    t  = np.arange(0, t_end, dt)
    D  = np.zeros(len(t))
    S  = np.zeros(len(t)); S[0] = S_rest
    TH = np.full(len(t), th_0)
    fires = []

    # precompute synaptic input at each timestep
    I_syn = np.zeros(len(t))
    for i, ti in enumerate(t):
        I_syn[i] = I_func(ti)
        if synapses:
            for w, spike_times in synapses:
                for ts in spike_times:
                    if ts <= ti:
                        I_syn[i] += w * np.exp(-(ti - ts) / tau_D)

    for i in range(1, len(t)):
        d = D[i-1]; s = S[i-1]; th = TH[i-1]
        dD  = (-d  + I_syn[i]) / tau_D
        dS  = (-(s - S_rest) + kappa * d - (th - th_0)) / tau_S
        dth = -(th - th_0) / tau_th
        D[i]  = d  + dD  * dt
        S[i]  = s  + dS  * dt
        TH[i] = th + dth * dt
        if S[i] >= 0:         # fire condition
            fires.append(t[i])
            S[i]  = S_reset
            TH[i] = th + dth * dt + dth   # theta jumps
    V = spike_waveform(t, fires)
    return t, D, S, TH, V, fires
\end{lstlisting}

\subsection{\texttt{dstc\_model.py}: The NMDA Gate and Calcium}

The NMDA gate is a single line. The full DSTC simulator adds the distal compartment, inhibition, and calcium:

\begin{lstlisting}[caption={NMDA gate and DSTC step (\texttt{dstc\_model.py})}, label={lst:dstc}]
def g_NMDA(V):
    """Voltage-gated NMDA conductance.
    Nearly 0 at rest (V=-1), ~1 when depolarized (V=2)."""
    return 1.0 / (1.0 + gamma_N * np.exp(-beta_N * V))

# Inside the integration loop:
nmda_in = g_NMDA(dd) * I_nmda_d(t[i])
inh_d   = I_inh_d(t[i]) * (dd - E_inh)   # shunting inhibition
inh_p   = I_inh_p(t[i]) * (dp - E_inh)

dDd = (-dd + I_syn_d[i] + nmda_in - inh_d) / tau_Dd
dDp = (-dp + I_syn_p[i] + kappa_dd * dd - inh_p) / tau_Dp
dS  = (-(s - S_rest) + kappa * dp - (th - th_0)) / tau_S
dth = -(th - th_0) / tau_th
dCa = -ca / tau_Ca

# ... update state ...

if S[i] >= 0:           # fire
    fires.append(t[i])
    S[i]  = S_reset
    TH[i] = th + dth * dt + dth
    Ca[i] = ca + dCa * dt + delta_Ca   # Ca jump
\end{lstlisting}

\subsection{\texttt{D\_compartment.py}: The Physical Circuit Simulator}

The D compartment is a capacitor and resistor. The simulator integrates Kirchhoff's current law directly, making it a direct correspondence between equation and component:

\begin{lstlisting}[caption={RC circuit simulator (\texttt{D\_compartment.py})}, label={lst:rc}]
# Circuit: tau_D = R_bleed * C_cap = 1kOhm * 10uF = 10ms
C_cap   = 10e-6    # 10 uF capacitor
R_bleed = 1e3      # 1 kOhm bleeder resistor (the leak)
R_syn   = 2e3      # 2 kOhm per synapse (weight = 1/R_syn)

def simulate_RC(input_sources, t_end=0.08, dt=0.0001):
    t = np.arange(0, t_end, dt)
    V = np.zeros(len(t))   # voltage across capacitor
    for i in range(1, len(t)):
        v    = V[i-1]
        I_leak = v / R_bleed      # current through bleeder (leak)
        I_in   = 0.0
        for amp, ts, te, rw in input_sources:
            if ts <= t[i-1] < te:
                I_in += (amp - v) / rw   # current through weight resistor
        dV   = (I_in - I_leak) / C_cap
        V[i] = v + dV * dt
    return t * 1000, V    # return time in ms
\end{lstlisting}

\noindent The equation $\tau_D \dot{V} = -V + \sum_i w_i I_i(t)$ is exactly what this circuit computes: \texttt{C\_cap} is $C$, \texttt{R\_bleed} is $R$ (so $\tau_D = RC$), and \texttt{R\_syn} sets the synaptic weight ($w_i = 1/R_w$).

% ── 6. Simulation Results ────────────────────────────────────────────────────
\section{Simulation Results}
\label{sec:results}

Each model file produces an 8-panel figure. We describe the key demonstrations.

\subsection{DST Behavioral Repertoire (\texttt{dst\_model.py})}

\textbf{Single spike (Panel 1).} A 2\,ms current injection at $t = 10$\,ms produces a clean action potential: $D$ charges, $S$ crosses threshold, $\theta$ jumps, $S$ resets below rest, $\theta$ recovers over $\sim$80\,ms.

\textbf{Refractory (Panel 3).} A second identical stimulus 4\,ms after the first is completely blocked ($\theta$ is still elevated). A third stimulus at $t = 28$\,ms fires successfully once $\theta$ has recovered.

\textbf{Temporal summation (Panel 4).} A single sub-threshold pulse fails to fire. Four identical pulses spaced 4\,ms apart (within the $\tau_D = 10$\,ms integration window) accumulate in $D$ and fire the neuron. The same pulses spaced 10\,ms apart do not.

\textbf{Spatial summation (Panel 5).} One synapse of weight 0.5 is sub-threshold. Three simultaneous synapses of the same weight fire the neuron. Six synapses produce a rapid burst.

\textbf{Spike-rate adaptation (Panel 6).} Under sustained current injection, each spike raises $\theta$, so the next spike requires $\theta$ to recover further before $S$ can reach threshold again. The inter-spike interval grows monotonically: the neuron adapts its firing rate to sustained input, a hallmark of biological neurons.

\subsection{DSTC Extensions (\texttt{dstc\_model.py})}

\textbf{NMDA gate (Panel 2).} $g_{\text{NMDA}}(V)$ is plotted across $V \in [-3, 4]$. At $V = -1$ (rest): $g = 0.027$. At $V = 0$ (threshold): $g = 0.25$. At $V = 2$ (depolarized): $g = 0.98$. The gate traverses three orders of magnitude across the physiological operating range.

\textbf{Coincidence detection (Panel 3).} NMDA input alone ($g_{\text{NMDA}}$ blocks it at rest) does not fire the neuron even at large amplitude. The same NMDA input combined with concurrent AMPA input (which depolarizes $D_d$, opening the gate) fires a burst of spikes. This demonstrates two-input coincidence detection in hardware.

\textbf{Distal-to-proximal delay (Panel 4).} A distal input drives $D_d$ (peak at $\sim$15\,ms), which then charges $D_p$ (peak delayed by several ms). The propagation delay through the dendritic cable is visible directly.

\textbf{Inhibitory shunting (Panel 5).} Proximal inhibition (blocking $D_p$) is more effective than distal inhibition (blocking $D_d$) at the same conductance. A stimulus that fires the neuron without inhibition is blocked by proximal inhibition and only weakened by distal inhibition.

\textbf{Calcium accumulation (Panel 6).} Under sustained input, $Ca$ builds up across spikes and decays slowly ($\tau_{Ca} = 200$\,ms). This staircasing calcium trace is the plasticity signal: a synapse consistently active before spikes will see rising $Ca$ and undergo LTP.

\subsection{D Compartment, Simulated Circuit Response (\texttt{D\_compartment.py})}

\noindent The figures in this subsection are produced by a circuit simulator, not measured on hardware. The component values are those specified in Section~\ref{sec:hardware}; the traces are what those values predict.

\textbf{Circuit (Panel 1).} The schematic: three synaptic inputs each through a 2\,k$\Omega$ weight resistor, converging on a 10\,$\mu$F capacitor with a 1\,k$\Omega$ bleeder to ground. Output is the capacitor voltage $V_D$.

\textbf{$\tau$ verification (Panel 2).} A single 5\,ms input pulse charges $V_D$. The decay from peak is measured; the time to reach $V_{\text{peak}} \cdot e^{-1}$ is $\tau_D = 10$\,ms, confirming $RC = 10$\,ms.

\textbf{Spatial summation (Panel 3).} One synapse peaks at $\sim$0.22\,V. Two simultaneous synapses at $\sim$0.39\,V. Three at $\sim$0.52\,V. The D-to-S handoff threshold is set at 0.45\,V: one synapse sub-threshold, two borderline, three fires.

\textbf{Temporal summation (Panel 4).} Two pulses 8\,ms apart (within $\tau_D$) stack: the second arrives before the first has decayed, accumulating on the residual charge. Two pulses 45\,ms apart fully decay between them.

% ── 7. Hardware Build Path ────────────────────────────────────────────────────
\section{Hardware Build Path}
\label{sec:hardware}

\subsection{D Compartment (Specified; Build Pending)}

\begin{center}
\small
\begin{tabular}{lll}
\toprule
\textbf{Component} & \textbf{Value} & \textbf{Function} \\
\midrule
Electrolytic capacitor $C$ & 10\,$\mu$F & Integrates charge (the memory) \\
Bleeder resistor $R$ & 1\,k$\Omega$ & Leak to ground ($\tau_D = RC = 10$\,ms) \\
Synaptic resistors $R_w$ & 1--10\,k$\Omega$ & Weight: $w_i = 1/R_{w,i}$ \\
\bottomrule
\end{tabular}
\end{center}

\noindent The bench for this measurement is implemented in \texttt{bench/} (firmware, capture, and an overlay that fits $\tau_D$ from the recorded trace and reports the capacitance it implies). No measured data is reported in this paper; the protocol below is what the bench executes.

\noindent\textbf{Test protocol.} (1) Apply a 2\,V pulse for 5\,ms through one synaptic resistor. Measure $V_D$ with an oscilloscope. Confirm exponential decay with $\tau = 10$\,ms. (2) Apply two simultaneous pulses; confirm $V_D$ peaks at $\sim 2\times$ the single-pulse value. (3) Apply two pulses 8\,ms apart; confirm $V_D$ from the second pulse builds on residual from the first. These three tests validate spatial summation, temporal summation, and the $\tau_D$ time constant.

\subsection{S Compartment (Next)}

The soma is a bistable threshold element with self-reset. Hardware: a Schmitt trigger. When $V_D$ crosses the upper threshold, the output switches high (spike); an RC feedback path charges a capacitor that drives the output back low after $\sim$1\,ms (reset). The adaptive threshold $\theta$ is an additional slow RC ($80\,\mu$F $\times$ 1\,k$\Omega$, $\tau = 80$\,ms) charged by the spike output.

\subsection{DSTC Full Circuit}

\begin{center}
\small
\begin{tabular}{llll}
\toprule
\textbf{Compartment} & \textbf{Hardware} & \textbf{$\tau$} & \textbf{Status} \\
\midrule
$D_d$ (distal)   & 15\,$\mu$F $\times$ 1\,k$\Omega$ RC + MOSFET gate & 15 ms & Designed \\
$D_p$ (proximal) &  5\,$\mu$F $\times$ 1\,k$\Omega$ RC                &  5 ms & Designed \\
$S$ (soma)       & Schmitt trigger + reset RC                          &  1 ms & Next \\
$\theta$ (threshold) & 80\,$\mu$F $\times$ 1\,k$\Omega$ slow RC       & 80 ms & Next \\
$Ca$ (calcium)   & 200\,$\mu$F $\times$ 1\,k$\Omega$ slow RC          & 200 ms & Planned \\
\bottomrule
\end{tabular}
\end{center}

\noindent The MOSFET gate for the NMDA nonlinearity uses the distal dendritic voltage $V_{D_d}$ as the gate control. At rest ($V_{D_d} \approx 0$\,V on the physical scale), the MOSFET is nearly pinched off (conductance low). As $V_{D_d}$ rises (AMPA depolarization), the gate opens and the NMDA current path is enabled. The physical sigmoid matches $g_{\text{NMDA}}(V)$ in the simulation.

% ── 8. Integration into SPECTRUM ─────────────────────────────────────────────
\section{Integration into SPECTRUM}
\label{sec:integration}

\subsection{DSTC as L1}

A network of DSTC neurons instantiates SPECTRUM's L1 (Perception layer). Each neuron exposes three signals to the common 1024-dimensional embedding space:

\begin{enumerate}
    \item \textbf{Spike rate}: spike count over a short integration window (rate-coded scalar).
    \item \textbf{Calcium level}: $Ca$ voltage, reflecting recent spike history and plasticity state.
    \item \textbf{Inhibitory control}: a scalar from L2 (working memory) that gates $g_{\text{inh},p}$, implementing top-down attention from the Hopfield layer.
\end{enumerate}

\noindent The NMDA coincidence detection means L1 is sensitive to \textit{correlated} inputs, not just strong inputs. The STDP via calcium means L1 learns online without backpropagation. Both properties are absent from existing clock-driven digital neuromorphic chips (Intel Loihi, IBM NorthPole), which implement simplified leaky integrate-and-fire (LIF) neurons without dendritic structure.

\subsection{ARBOR as L6}

L6 (Language layer) is instantiated as the ARBOR architecture \cite{lin2026arbor}: a \textbf{Big Brain} generalist model handles routing and synthesis; a registry of \textbf{Long Brain} specialists (each trained exclusively on one domain) provides depth. The Big Brain receives the embedding from L4, identifies required domains, routes to Long Brains, and synthesizes responses.

ARBOR and SPECTRUM are orthogonal: SPECTRUM defines \textit{what physical substrates} compute each function; ARBOR defines \textit{how domain knowledge is organized} at the output stage. A SPECTRUM system reasoning about quantum chemistry uses L5 for hypothesis search, L4 for continuous optimization, and the chemistry Long Brain at L6 for domain-precise output. The Big Brain synthesizes across all three.

\subsection{FLUID as Dynamic Topology}

The FLUID architecture \cite{lin2026fluid}, which modifies neural topology via Hebbian co-activation during inference, runs within each ARBOR Long Brain. This provides morphological computation at the software level inside L6---topology-as-computation without requiring hardware-level reconfigurable photonic meshes.

% ── 9. Related Work ──────────────────────────────────────────────────────────
\section{Related Work}
\label{sec:related}

\textbf{Hodgkin-Huxley} \cite{hodgkin1952} is the gold standard biological neuron model. DST deliberately departs from it: HH is calibrated to squid axon data; DST is designed for computation. HH requires four variables and six conductance functions; DST requires three variables and no ion channel fits.

\textbf{Leaky integrate-and-fire (LIF)} \cite{lapicque1907} is the industry standard for neuromorphic chips. LIF has one variable; it cannot produce spike-rate adaptation, NMDA coincidence detection, or dendritic integration. DST adds these at the cost of two additional variables.

\textbf{Izhikevich} \cite{izhikevich2003,izhikevich2004} is the closest prior work and sets the bar DST must clear. It reaches spike-rate adaptation and a wide firing repertoire with \textit{two} state variables and four parameters, and it was written for the reason given in Section~\ref{sec:dst}---that Hodgkin-Huxley is too expensive to put on a chip. We state plainly that DST does not win on economy: three variables and eight parameters is more, not less. Table~\ref{tab:baselines} reports what DST does buy for that cost.

\textbf{Adaptive exponential integrate-and-fire (AdEx)} \cite{brette2005} likewise obtains adaptation with two variables, via an exponential spike-initiation term and a coupled adaptation current.

\begin{table}[htbp]
\centering
\small
\caption{DST against published spiking-neuron models. All models integrated at
$dt = 0.01$\,ms over 300\,ms; each driven at multiples of its \textit{own}
rheobase, since the four carry different current units. Behaviours are read
from each model's output rather than asserted. ``Max stable $dt$'' is the
coarsest timestep that still reproduces the spike count obtained at
$dt = 0.01$\,ms. Runtime is wall-clock for one 300\,ms trace, same host, same
$dt$. Reproduced by \texttt{bench/baselines.py}.}
\label{tab:baselines}
\begin{tabular}{lccccccc}
\toprule
\textbf{Model} & \textbf{Vars} & \textbf{Params} & \textbf{Refractory}
  & \textbf{Temporal sum.} & \textbf{Adaptation}
  & \textbf{Max $dt$ (ms)} & \textbf{Runtime (ms)} \\
\midrule
\textbf{DST (ours)} & 3 & 8 & \textbf{yes} & yes & yes & \textbf{2.0} & \textbf{12.7} \\
LIF                 & 1 & 5 & no  & yes & no  & 1.0 & 6.8 \\
Izhikevich          & 2 & 4 & no  & yes & yes & 1.0 & 23.2 \\
AdEx                & 2 & 9 & no  & yes & yes & 0.5 & 68.1 \\
\bottomrule
\end{tabular}
\end{table}

Two results in Table~\ref{tab:baselines} bear on the design. First, DST is the
only model tested that holds an \textit{absolute} refractory period: under a
pulse pair at 20$\times$ rheobase the others recover through a slow current and
can be driven to fire again immediately, whereas the threshold jump $\Delta\theta$
is a hard veto. Second, and more consequential for fabrication, DST tolerates a
timestep twice Izhikevich's and four times AdEx's, and at matched $dt$ runs
$1.8\times$ faster than Izhikevich and $5.4\times$ faster than AdEx despite the
extra state variable. Both competitors pay for a nonlinearity---a quadratic term
requiring sub-stepping, an exponential---while DST's dynamics are entirely
linear.

That linearity is the argument, not the variable count. A linear system maps
onto an RC network, which is why the D compartment in Section~\ref{sec:hardware}
is a capacitor and a resistor; neither a quadratic nor an exponential term has
any comparable two-component realisation. The point-neuron comparison also
cannot score the DSTC extensions at all, since a point neuron has no morphology
on which to place a two-compartment dendrite, a voltage-gated NMDA conductance,
or a proximal/distal inhibition asymmetry.

\textbf{Intel Loihi} \cite{davies2021} and \textbf{IBM NorthPole} \cite{ibm2023} are state-of-the-art neuromorphic chips using LIF neurons. Neither integrates with transformer language models. SPECTRUM proposes this integration with DSTC at L1.

\textbf{Multi-compartment models} \cite{mainen1996} extend HH with spatial dendritic structure. DSTC takes a minimal two-compartment approach (distal + proximal) sufficient to demonstrate NMDA coincidence detection and shunting inhibition without the parameter burden of full compartmental models.

\textbf{STDP} \cite{bi1998} is well-established experimentally. The calcium-based STDP in DSTC follows the framework of \cite{shouval2002}, using $Ca$ as the plasticity induction variable rather than directly tracking spike timing.

% ── 10. Conclusion ────────────────────────────────────────────────────────────
\section{Conclusion}
\label{sec:conclusion}

The argument for SPECTRUM is simple: intelligence needs different physical machinery for different cognitive functions, and building it all on one substrate is leaving capability and energy efficiency on the table. The transformer is not wrong; it is mispositioned. It belongs at the output, not the core.

The DST neuron shows that a neuromorphic unit can be derived from first principles---three equations, three variables, eight parameters---and produce the behavioral repertoire of a spiking neuron without biological fitting. That repertoire alone is not the contribution: Izhikevich reaches it with two variables and four parameters \cite{izhikevich2003}. What DST adds, per Table~\ref{tab:baselines}, is an absolute refractory period the two-variable models lack, and fully linear dynamics---which cost less to integrate and, unlike a quadratic or exponential term, correspond directly to passive components. DSTC extends this with NMDA coincidence detection, two-compartment dendritic integration, inhibitory shunting, and calcium-based online learning. These are not biological luxuries; they are computational capabilities (coincidence detection, online STDP, top-down inhibitory control) that make L1 more powerful than a LIF array.

The D compartment circuit---10\,$\mu$F, 1\,k$\Omega$, a few weight resistors---is the first physical piece. It is buildable today on a breadboard, and has not yet been built; the measurement bench is written and awaiting the parts. The progression from there to the full DSTC chip is specified: S compartment (Schmitt trigger), $\theta$ compartment (slow RC), calcium compartment (slow RC on spike output), MOSFET NMDA gate. Each step is a standard analog component.

The Python code (\texttt{dst\_model.py}, \texttt{dstc\_model.py}, \texttt{D\_compartment.py}) is not a simulation to be discarded---it is the specification. Every equation in the code corresponds directly to a circuit element. The software and hardware are the same system at different levels of abstraction.

SPECTRUM needs six layers. Three now have concrete instantiations: DSTC at L1, analog at L4, ARBOR at L6. The integration is the research agenda.

% ── References ────────────────────────────────────────────────────────────────
\begin{thebibliography}{99}

\bibitem{vaswani2017} Vaswani, A., et al.\ (2017). Attention is all you need. \textit{NeurIPS 2017}.

\bibitem{hodgkin1952} Hodgkin, A.L., \& Huxley, A.F.\ (1952). A quantitative description of membrane current and its application to conduction and excitation in nerve. \textit{Journal of Physiology}, 117(4), 500--544.

\bibitem{izhikevich2003} Izhikevich, E.M.\ (2003). Simple model of spiking neurons. \textit{IEEE Transactions on Neural Networks}, 14(6), 1569--1572.

\bibitem{izhikevich2004} Izhikevich, E.M.\ (2004). Which model to use for cortical spiking neurons? \textit{IEEE Transactions on Neural Networks}, 15(5), 1063--1070.

\bibitem{brette2005} Brette, R., \& Gerstner, W.\ (2005). Adaptive exponential integrate-and-fire model as an effective description of neuronal activity. \textit{Journal of Neurophysiology}, 94(5), 3637--3642.

\bibitem{lapicque1907} Lapicque, L.\ (1907). Recherches quantitatives sur l'excitation \'{e}lectrique des nerfs trait\'{e}e comme une polarisation. \textit{J.\ Physiol.\ Pathol.\ Gen.}, 9, 620--635.

\bibitem{davies2021} Davies, M., et al.\ (2021). Advancing neuromorphic computing with Loihi. \textit{Proceedings of the IEEE}, 109(5), 911--934.

\bibitem{ibm2023} IBM Research (2023). NorthPole: An architecture for neural network inference with a 12nm chip. \textit{Science}, 382(6668).

\bibitem{mainen1996} Mainen, Z.F., \& Sejnowski, T.J.\ (1996). Influence of dendritic structure on firing pattern in model neocortical neurons. \textit{Nature}, 382, 363--366.

\bibitem{bi1998} Bi, G., \& Poo, M.\ (1998). Synaptic modifications in cultured hippocampal neurons. \textit{Journal of Neuroscience}, 18(24), 10464--10472.

\bibitem{shouval2002} Shouval, H.Z., et al.\ (2002). A unified model of NMDA receptor-dependent bidirectional synaptic plasticity. \textit{PNAS}, 99(16), 10831--10836.

\bibitem{lin2026arbor} Lin, W.Z.\ (2026). ARBOR: Adaptive Routing with Branch-Organized Reasoning. \textit{arXiv preprint}.

\bibitem{lin2026fluid} Lin, W.Z.\ (2026). FLUID: A self-structuring architecture for open-ended artificial intelligence. \textit{arXiv preprint}.

\bibitem{lin2026spectrum} Lin, W.Z.\ (2026). SPECTRUM: A Multi-Substrate Hybrid Architecture for General-Purpose Artificial Intelligence. \textit{arXiv preprint}.

\end{thebibliography}

\end{document}
